% Generated draft: requires mathematical and source review.
\documentclass[11pt]{article}
\usepackage[margin=25mm]{geometry}
\usepackage{fontspec}
\setmainfont{Latin Modern Roman}
\usepackage{amsmath,amssymb,mathtools,unicode-math}
\setmathfont{Latin Modern Math}
\usepackage{enumitem,needspace,xurl,hyperref}
\setlength{\parindent}{0pt}
\setlength{\parskip}{5pt}
\setlength{\emergencystretch}{3em}
\newcommand{\N}{\mathbb{N}}
\newcommand{\Z}{\mathbb{Z}}
\newcommand{\Q}{\mathbb{Q}}
\newcommand{\R}{\mathbb{R}}
\newcommand{\C}{\mathbb{C}}
\newcommand{\F}{\mathbb{F}}
\newcommand{\A}{\mathbb{A}}
\newcommand{\PP}{\mathbb{P}}
\newcommand{\E}{\mathbb{E}}
\newcommand{\eps}{\varepsilon}
\newcommand{\dd}{\,\mathrm{d}}
\begin{document}
\section*{UnsolvedMath ID: 30005564}
\textbf{Source problem code:} catalogue entry
\subsection{Definitions and mathematical statement}
For the characteristic-zero formulation explicitly mentioned in the source, let \(K\) be an algebraically closed field of characteristic zero. For a smooth, connected, nondegenerate curve \(C\subset\PP^r_K\) of degree \(d\) and genus \(g\), with \(r\geq2\), define its normal bundle by
\[
N_{C/\PP^r}=\mathcal{H}om_{\mathcal O_C}(\mathcal I_C/\mathcal I_C^2,\mathcal O_C).
\]
It has rank \(r-1\) and degree \((r+1)d+2g-2\). For a vector bundle \(E\) of positive rank on \(C\), write \(\mu(E)=\deg(E)/\operatorname{rk}(E)\). The bundle is slope stable, respectively slope semistable, if every nonzero proper saturated subbundle \(F\subset E\) of smaller rank satisfies
\[
\mu(F)<\mu(E),\qquad\text{respectively}\qquad\mu(F)\leq\mu(E).
\]
Thus stability of the normal bundle means
\[
\frac{\deg(F)}{\operatorname{rk}(F)}<\frac{(r+1)d+2g-2}{r-1}
\]
for every such \(F\subset N_{C/\PP^r}\).

The Brill--Noether number is
\[
\rho(d,g,r)=g-(r+1)(g-d+r)=(r+1)d-rg-r(r+1).
\]
The supplied excerpt does not define its term ``BN triple'' or specify all admissible values of \(d,g,r\). The usual interpretation involves \(\rho(d,g,r)\geq0\) and a general Brill--Noether curve, in a parameter-space component dominating the moduli space of curves of genus \(g\). The numerical inequality alone is not adopted here as a complete definition of the source's domain: the allowed ambient dimensions, embedding hypotheses, and intended component require confirmation. In what follows, \(\mathcal T_{\mathrm{BN}}\) denotes precisely the set of admissible BN triples intended by the source, restricted to \(g\geq2\). Subject to that missing specification, call a triple stable when the normal bundle of a general BN curve of that triple is slope stable. Here ``general'' means that the property holds on a nonempty Zariski open subset of the intended irreducible parameter-space component.

\textbf{Conjectural statement, with the domain qualification above.} The characteristic-zero version asks for a finite set \(E\subset\mathcal T_{\mathrm{BN}}\) such that
\[
\forall(d,g,r)\in\mathcal T_{\mathrm{BN}}\setminus E,
\qquad N_{C/\PP^r}\text{ is stable for a general BN curve }C.
\]
Equivalently, the set of triples for which the general normal bundle is not stable should be finite, with finiteness taken over all ambient dimensions together. The stronger task is to determine this exceptional set. The original excerpt more specifically proposes that every admissible triple outside those in its Examples 1 and 2 is stable, at least in characteristic zero. If \(E_{\mathrm{ex}}\) denotes the triples in those examples, this asserts that the exceptional set is contained in \(E_{\mathrm{ex}}\); the examples themselves were not supplied, so neither their triples nor the claimed exceptional behavior can be specified here. The first finiteness assertion in the original excerpt has no explicit characteristic qualification; its intended scope outside characteristic zero remains unresolved.

\textbf{Supplied evidence and limitations.} The excerpt attributes the conjecture to [CLV22, Conjecture 1.1]. It also reports sufficient conditions for semistability of a general BN curve's normal bundle in [CG23, Theorem 1.1], including a genus-one case, and sufficient conditions for stability when \(g\geq2\) in [CG23, Theorem 1.2]. Its displayed genus and degree bounds, including the definition of the auxiliary parameter \(u\), are corrupted in the supplied transcription and cannot be stated reliably. The genus-one evidence lies outside the conjecture's \(g\geq2\) domain.

One legible auxiliary definition in that evidence is that \(b_2(r)\) is the least integer \(B\) such that every integer \(d\geq B\) admits positive integers \(d_1,d_2\) satisfying
\[
d=d_1+d_2,\qquad d_1,d_2\geq r+1,\qquad\gcd(r-1,2d_1+1)=1.
\]
The excerpt reports
\[
b_2(r)\leq2r+\frac{p-1}{2},\qquad
b_2(r)=2r+O(\log r),\qquad
b_2(r)\leq\frac{5r-3}{2},
\]
where \(p\geq5\) is the smallest prime not dividing \(r-1\), and additionally reports \(b_2(r)\leq2.01r+2.015\) for \(r\geq1636\). These are retained as supplied background claims, not independently verified assertions: the excerpt does not give a reliable complete range of validity for all these bounds, and its surrounding formulas are damaged.

The supplied dated literature assessment says that stability outside finitely many triples has been proved for each fixed \(r\), with enlarged stable and semistable ranges. Finiteness separately for each \(r\) does not supply the conjectured finiteness over all triples. This assessment and the supplied \texttt{partially\_solved} label are curation claims, not mathematical conclusions verified by this draft.
\subsection{Short English statement}
For a general Brill--Noether curve of genus at least two, conjecturally its normal bundle is slope stable for all but finitely many degree--genus--dimension triples. A stronger characteristic-zero conjecture confines the exceptions to two examples absent from the supplied record. The precise admissible triples and several corrupted background bounds require source review. The supplied claim of finiteness for each fixed dimension is weaker than finiteness across all dimensions.
\subsection{Sources}
\begin{itemize}
\item Local UnsolvedMath registry, \texttt{data/unsolvedmath/problems.json}, permanent ID \(30005564\), title ``Stability of Brill--Noether Triples,'' source code \texttt{catalogue entry}; the project also identifies \texttt{catalog/problems.json} as a supplementary registry. This draft uses the supplied complete record without independently inspecting those files.
\item Supplied primary-report citation: \emph{Recent Trends in Algebraic Geometry} (2023), Oberwolfach Reports, DOI \texttt{10.4171/owr/2023/27}, pp.\ 1485--1542. \url{https://doi.org/10.4171/owr/2023/27}. The record attributes the conjecture to [CLV22, Conjecture 1.1] and supporting results to [CG23, Theorems 1.1 and 1.2]; complete bibliography entries for these labels are absent. No individual proposer is supplied; the proposed year is 2023.
\item Supplied literature reference: Coskun--Larson--Vogt, \emph{Stability of natural bundles on curves}. \url{https://arxiv.org/abs/2602.03944}.
\item Supplied literature reference: Coskun--Larson--Vogt, \emph{Stability of normal bundles of curves}. The supplied link is the workshop-report DOI: \url{https://doi.org/10.4171/owr/2023/27}.
\item Literature assessment dated exactly \texttt{2026-08-22}: supplied dated curation, not independently verified. It labels the problem partially solved and says that the conjecture remains open, that stability is known outside finitely many triples for each fixed \(r\), and that known stable and semistable ranges have expanded. The background expressly requests independent verification of time-sensitive status and cited claims.
\item The registry labels the statement \texttt{corrected\_verified} and says it was reconstructed and checked against the official Oberwolfach report. These are supplied provenance claims; this draft does not claim to have visited or verified any source.
\end{itemize}
\end{document}
